\documentclass[degree=doctor,bibtype=numeric,degreetype=academic]{tongjithesis}
% 选项：
%   degree=[master|doctor], 							% 必选
%   bibtype=[numeric|authoryear], 						% 可选，数字式引用|作者-年份引用，默认为数字式（上标）引用
%   degreetype=[academic|profession|equaleducation],  	% 可选, 学术型|专业型|同等学力，默认为学术型
% 	electronic,                                 		% 可选, 电子版，(打印时删除)
%   secret,                                     		% 可选，是否保密，基本不用
%   pifootnote,                                 		% 可选，默认已打开
%   romantitle                                  		% 可选，默认已打开
%   注：默认已打开的选项可以使用arialtitle=false的形式关闭。

% 所有其它可能用到的包都统一放到这里了，可以根据自己的实际添加或者删除。

\usepackage{tongjiutils}
\usepackage{import}
\usepackage{algorithm}
\usepackage{algorithmic}
% \usepackage{graphicx}
% \usepackage{subfigure}
% \usepackage{amsmath} %equation
% \usepackage{mathabx} %rightrightharpoons




\newcommand{\para}[1]{\smallskip\noindent {\bf #1}}
\newcommand{\exampledp}{\texttt{ExampleDP}}
\newcommand{\rightrightharpoons}{{\Longrightarrow}}   % TODO: fix it
\newcommand{\codeword}[1]{\texttt{\small{#1}}}

\newcommand{\rulesep}{\unskip\ \vrule\ }
\newcommand \bnfdef  {\mathrel{\color{blue}{::=}}}
\newcommand \bnfalt  {\mathrel{\color{blue}{|}}}

\newcommand \keyfont {\mathsf}
\newcommand \ONPKT {\keyfont{onPacket(pkt)}}
% \newcommand \WHILE  {\keyfont{while}}
\newcommand \IFSTR     {\keyfont{if}}
\newcommand \THEN     {\keyfont{then}}
% \newcommand \SKIP   {\keyfont{skip}}
\newcommand \RETURNSTR   {\keyfont{return}}
% \newcommand \ADD   {\keyfont{add}}
% \newcommand \CONTAINS   {\keyfont{contains}}
% \newcommand \GET   {\keyfont{get}}
% \newcommand \PUT   {\keyfont{put}}
% \newcommand \GLOBAL   {\keyfont{global}}
% \newcommand \forl   {\keyfont{for}}
% \newcommand \while[2]   {\WHILE\ {#1}\ {#2}}
\newcommand \ifelse[3]  {\IFSTR\ {#1}:{#2}\ \THEN:{#3}}
\newcommand \return[1]   { \RETURNSTR\ {#1}}
% \newcommand \CYCLIC  {\keyfont{loop}}
% \newcommand \cyclic[1]   { \CYCLIC\ {#1}}

% \newcommand \JUMP  {\keyfont{jump}}
% \newcommand \CJUMP  {\keyfont{cjump}}
% \newcommand \LBL  {\keyfont{label}}

% \newcommand \jump[1]   { \JUMP\ {#1}}
% \newcommand \cjump[3]   { \CJUMP\ {#1}\ {#2}\ {#3}}
% \newcommand \lbl[1]   { \LBL\ {#1}}

\newcommand \tx {\mathtt{x}}
\newcommand \ty {\mathtt{y}}
\newcommand \tz {\mathtt{z}}
\newcommand \tg {\mathtt{g}}



\newtheorem{corrollary}{推论}

%参考文献更新使用biblatex包, 使用gb7714-2015标准, 具体参数设置可在cls文件中搜索biblatex进行了解
%加入bib文件(老版本文件依然能够使用)
\addbibresource{intro-cn/all-all.bib}   %


\begin{document}

% 定义所有的eps文件在 figures 子目录下
% \graphicspath{{figures/}}


%%% 封面部分
\frontmatter
\tongjisetup{
  %******************************
  % 注意：
  %   1. 配置里面不要出现空行
  %   2. 不需要的配置信息可以删除
  %******************************
  %
  %=====
  % 秘级
  %=====
  secretlevel={保密},
  secretyear={2},
  %
  %=========
  % 中文信息
  %=========
  % 题目过长可以换行（推荐手动加入换行符，这样就可以控制换行的地方啦）。
  ctitle={软件定义网络高级编程模型研究},
  cheadingtitle={软件定义网络高级编程模型研究},    %用于页眉的标题，不要换行 
  studentnumber={XXXXXXX},
  cdepartment={电子与信息工程学院},
  cmajorfirst={计算机科学与技术},
  cmajorsecond={计算机软件与理论},
  cauthor={XXX}, 
  csupervisor={XXX 教授}, 
  % 如果没有副指导老师或者校外指导老师，把{}中内容留空即可，或者直接注释掉。
  % cassosupervisor={裴刚 教授~（校外）}, % 副指导老师
  % 日期自动使用当前时间，若需手动指定，按如下方式修改：
  % cdate={\zhdigits{2018}年\zhnumber{11}月},
  % 没有基金的话就注释掉吧。
  % cfunds={（本论文由我要努力想办法撑到两行的著名国家杰出青年基金 (No.123456789) 支持）},
  %
  %=========
  % 英文信息
  %=========
  etitle={Software Defined Network\\ High-Level Programming Model}, 
  edepartment={College of Electronics and Information Engineering},
  emajorfirst={Computer Science and Technology},
  emajorsecond={Computer Software and Theory},
  % 日期自动使用当前时间，若需手动指定，按如下方式修改：
  % edate={November,\ 2018},
  % efunds={(Supported by the Natural Science Foundation of China for\\ Distinguished Young Scholars, Grant No.123456789)},    
  eauthor={XXX},
  esupervisor={Prof. XXX},
  % eassosupervisor={Prof. Gang Pei (XiaoWai)}
  }

% 定义中英文摘要和关键字
\begin{cabstract}  

软件定义网络（SDN）的一个重要能力是高级编程。通过SDN高级编程，网络管理者可以实现对网络中数据流的灵活控制，而不必关心底层网络设备的具体架构。然而，一方面，随着SDN的普及，新应用的不断产生，人们对SDN高级编程的要求也不断提高；另一方面，随着网络设备技术的迅速发展，底层网络设备在不断变得高效、灵活的同时，其架构也变得越来越复杂。因此，如何将高级SDN程序高效地转化为底层网络设备的低级配置，是SDN高级编程模型研究始终需要考虑的问题。

对于单交换机数据通路的情况，目前有通过面向低级配置的接口生成具有固定结构（如OF-DPA）数据通路配置的相关工作，也有将高级程序转化为具有可定制结构（如RMT）数据通路配置的相关工作，但是没有将高级程序转化为具有固定结构数据通路配置的相关工作。要想利用那些高效的但具有固定结构数据通路，而且不想关心其复杂的实现架构，则将高级程序转化为具有固定结构数据通路配置的工作非常重要。而在进行转化前，一个重要问题则是，一个高级SDN程序是否可以被具有固定结构数据通路所实现。该问题的复杂地方在于高级SDN程序和底层数据通路之间结构差异性非常大，而且各自都具有丰富的多样性，因此我们需要一个系统的方法去解决该问题。针对该问题，我们提出一个将高级SDN程序和底层数据通路统一的特征空间，通过特征函数，高级SDN程序和数据通路都可以映射到该空间上的点，最后通过定义空间上点的比较函数，来判断该高级SDN程序是否可以被该数据通路实现。除此之外，当高级SDN程序具有循环结构时，如何将该程序高效地实现在可定制结构的数据通路上，也没有有效的工作。针对该问题，我们提出了重复软件流水线模型。该模型的优势在于可以对具有动态循环条件的循环结构进行转换，以及在转换后的结构上计算最佳数据通路时显示出较高的效率。

对于全网络数据平面的情况，大量的工作仍然是以响应式的模式（即不会主动生成配置，而等待数据包传递给SDN控制器）对网络进行管理，因此效率并不高。SNAP实现了从高级SDN程序经过主动编译转化为全网络且支持有状态的数据平面配置。然而，其采用One-Big-Switch模型，导致用户无法让数据包按指定的路径转发。因此，我们这里考虑的是一个非常灵活的高级SDN程序，即程序的输出为网络中的路径而非目的主机。并且，当网络存在具有固定功能的网络节点（如中间盒）时，我们的编程模型支持在高级SDN程序中调用相应网络节点功能对数据包进行处理，并根据调用返回的结果，用户可以选择网络中不同路径对数据包进行转发。对于这样的高级SDN程序，由于调用相应网络节点处理数据包时，数据包需要转发到该节点，而这可能会破坏程序中对数据包规定的转发路径的约束。因此，我们进一步考虑了其程序正确性问题，并提出系统路径约束来解决该问题。除此之外，针对不同网络场景，如软件定义联合网络的低时延要求，以及车联网中车的移动性要求，如何优化数据平面配置是一个挑战。为此，针对软件定义联合网络的场景，我们设计高级编程系统，并基于共享本地状态的方法，提高系统性能；针对车联网的场景，我们提出软件定义车联网的架构以及编程框架，并考虑车联网中车移动性的特点，提出优化的规则下发方法，减少生成的规则数量。

\end{cabstract}

\ckeywords{软件定义网络，可编程数据平面，高级编程模型}

\begin{eabstract}

A key capability of Software Defined Networking (SDN) is high-level programming. By using SDN high-level programming language, network managers can achieve flexible control of data flows in the network without the knowledge of the specific architecture of the underlying network devices. However, on the one hand, as SDN becoming more and more popular, new applications appears continuously. People's requirements for SDN high-level programming languages ​​are constantly improving; on the other hand, with the rapid development of network device technology, the underlying packet handling is becoming efficient and also flexible. However, at the same time, the data path architecture is becoming more and more complex. Therefore, how to efficiently convert high-level SDN programs into low-level configurations of underlying network devices is an issue that must be considered in SDN high-level programming model.



For the case of single switch data path, there are studies focusing on generating fixed-structure (such as OF-DPA) data path configuration through a low-level configuration interface, and also there are studies focusing on transforming a high-level program into a customizable structure (such as RMT) data path configuration. However, there is little work transforming a high-level program into a fixed-structured data path configuration To take advantage of efficient but fixed-structure data paths and do not want to care about their complex implementation architectures, it is important to turn high-level programs into the fixed-structured data path configurations. Then, an important issue before the transformation is whether a high-level SDN program can be implemented with a fixed-structure data path. This problem is not easy as the structural differences between the high-level SDN program and the underlying data path. So we need a systematic approach to solve the problem. To solve the problem, we propose a characteristic space that unifies the high-level SDN program and the underlying data path. Through the characteristic function, the high-level SDN program and the data path can be mapped to the points on the space. And by defining the comparison function of the points on the space, we can determine whether the high-level SDN program can be implemented by the data path. In addition, when the high-level SDN program has a loop structure, how to efficiently implement the program on the data path of the customizable structure is another problem. For this problem, we propose a repeated software pipeline model. The advantage of this model is that it can convert loop structures with dynamic loop conditions. And it is very efficient when calculating the optimal data path on the converted structure.




For the case of multiple data paths, a lot of studies are still in the reactive mode (that is, it does not actively generate the configuration, but waits for the packet to be delivered to the SDN controller) to manage the network, so the efficiency is not high. SNAP implements proactive compilation from a high-level SDN program to configurations of multiple stateful switches. However, it uses the One-Big-Switch model in which users cannot forward packets to a specified path in the network. What we are considering is a very flexible high-level SDN program where the user can calculate the path in the network as a return of the program. Moreover, when the network has a fixed-function node (such as a middlebox), our programming model supports invoking the corresponding node's function in the high-level SDN program to process packets, and based on the result returned by the invocation, the user can select different paths for packets to forward. For such a high-level SDN program, packets need to be forwarded to the node for the processing, which may break the constraints on the forwarding path specified in the program. Therefore, we further consider the correctness of the program and propose system path constraints to solve the problem. In addition, data plane configuration optimization is a challenge for different network scenarios, such as the low latency requirements of software defined coalition network and the mobility requirements of vehicles in the Internet of Vehicles. To this end, for the scenario of software defined coalition network, we design a high level programming system and improve the performance based on the method of sharing local state. For the scenario of Internet of Vehicles, we propose the Software Defined Internet of Vehicles architecture. For mobility in the network, the improved rule installation is proposed, by which the number of generated rules is reduced.

\end{eabstract}

\ekeywords{SDN, Programmable Datapath, High-Level Programming Model}
\makecover


% 目录
\tableofcontents
% 符号对照表
% \begin{denotation}
% \item[GNU] GNU's Not Unix /'gnu:/
\item[GFDL] GNU Free Documentation License
\item[GPL] GNU General Public License
\item[FSF] Free Software Foundation
\item[SMP] 对称多处理
\item[API] 应用程序编程接口
\item[$E$] 能量
\item[$m$] 质量
\item[$c$] 光速
\item[$P$] 概率
\item[$T$] 时间
\item[$v$] 速度

% \end{denotation}

%%% 以下索引按需要选择
% 插图索引
\listoffigures
% 表格索引
\listoftables
% 公式索引
% \listofequations

%%% 正文 
\mainmatter
\import{./intro-cn/}{intro-chap.tex}
\import{./cap-cn/}{cap-chap.tex}
\import{./loop-cn/}{loop-chap.tex}
\import{./global-cn/}{global-chap.tex}
\import{./stateful-cn/}{stateful-chap.tex}
\import{./sdiv-cn/}{sdiv-chap.tex}





%%% 其它部分
\backmatter

% 致谢
% \begin{acknowledgement}
% 衷心感谢导师 xxx 教授和物理系 xxx 副教授对本人的精心指导。他们的言传身教将使
我终生受益。

在美国麻省理工学院化学系进行九个月的合作研究期间，承蒙 xxx 教授热心指导与帮助，不
胜感激。感谢 xx 实验室主任 xx 教授，以及实验室全体老师和同学们的热情帮助和支
持！本课题承蒙国家自然科学基金资助，特此致谢。

感谢 \tongjithesis{}，它的存在让我的论文写作轻松自在了许多，让我的论文格式规整漂亮了
许多。
% \end{acknowledgement}
\clearpage
\newpage
\mbox{}
\newpage

% 参考文献
\printTJbibliography


% 附录
\begin{appendix}
% \chapter{外文资料原文}
\label{cha:engorg}
As one of the most widely used techniques in operations research, {\em
  mathematical programming} is defined as a means of maximizing a quantity known
as {\em objective function}, subject to a set of constraints represented by
equations and inequalities. Some known subtopics of mathematical programming are
linear programming, nonlinear programming, multiobjective programming, goal
programming, dynamic programming, and multilevel programming$^{[1]}$.

It is impossible to cover in a single chapter every concept of mathematical
programming. This chapter introduces only the basic concepts and techniques of
mathematical programming such that readers gain an understanding of them
throughout the book$^{[2,3]}$.


\section{Single-Objective Programming}
The general form of single-objective programming (SOP) is written
as follows,
\begin{equation}\tag*{(123)} % 如果附录中的公式不想让它出现在公式索引中，那就请
                             % 用 \tag*{xxxx}
\left\{\begin{array}{l}
\max \,\,f(x)\\[0.1 cm]
\mbox{subject to:} \\ [0.1 cm]
\qquad g_j(x)\le 0,\quad j=1,2,\cdots,p
\end{array}\right.
\end{equation}
which maximizes a real-valued function $f$ of
$x=(x_1,x_2,\cdots,x_n)$ subject to a set of constraints.

\newtheorem{mpdef}{Definition}[chapter]
\begin{mpdef}
In SOP, we call $x$ a decision vector, and
$x_1,x_2,\cdots,x_n$ decision variables. The function
$f$ is called the objective function. The set
\begin{equation}\tag*{(456)} % 这里同理，其它不再一一指定。
S=\left\{x\in\Re^n\bigm|g_j(x)\le 0,\,j=1,2,\cdots,p\right\}
\end{equation}
is called the feasible set. An element $x$ in $S$ is called a
feasible solution.
\end{mpdef}

\newtheorem{mpdefop}[mpdef]{Definition}
\begin{mpdefop}
A feasible solution $x^*$ is called the optimal
solution of SOP if and only if
\begin{equation}
f(x^*)\ge f(x)
\end{equation}
for any feasible solution $x$.
\end{mpdefop}

One of the outstanding contributions to mathematical programming was known as
the Kuhn-Tucker conditions\ref{eq:ktc}. In order to introduce them, let us give
some definitions. An inequality constraint $g_j(x)\le 0$ is said to be active at
a point $x^*$ if $g_j(x^*)=0$. A point $x^*$ satisfying $g_j(x^*)\le 0$ is said
to be regular if the gradient vectors $\nabla g_j(x)$ of all active constraints
are linearly independent.

Let $x^*$ be a regular point of the constraints of SOP and assume that all the
functions $f(x)$ and $g_j(x),j=1,2,\cdots,p$ are differentiable. If $x^*$ is a
local optimal solution, then there exist Lagrange multipliers
$\lambda_j,j=1,2,\cdots,p$ such that the following Kuhn-Tucker conditions hold,
\begin{equation}
\label{eq:ktc}
\left\{\begin{array}{l}
    \nabla f(x^*)-\sum\limits_{j=1}^p\lambda_j\nabla g_j(x^*)=0\\[0.3cm]
    \lambda_jg_j(x^*)=0,\quad j=1,2,\cdots,p\\[0.2cm]
    \lambda_j\ge 0,\quad j=1,2,\cdots,p.
\end{array}\right.
\end{equation}
If all the functions $f(x)$ and $g_j(x),j=1,2,\cdots,p$ are convex and
differentiable, and the point $x^*$ satisfies the Kuhn-Tucker conditions
(\ref{eq:ktc}), then it has been proved that the point $x^*$ is a global optimal
solution of SOP.

\subsection{Linear Programming}
\label{sec:lp}

If the functions $f(x),g_j(x),j=1,2,\cdots,p$ are all linear, then SOP is called
a {\em linear programming}.

The feasible set of linear is always convex. A point $x$ is called an extreme
point of convex set $S$ if $x\in S$ and $x$ cannot be expressed as a convex
combination of two points in $S$. It has been shown that the optimal solution to
linear programming corresponds to an extreme point of its feasible set provided
that the feasible set $S$ is bounded. This fact is the basis of the {\em simplex
  algorithm} which was developed by Dantzig as a very efficient method for
solving linear programming.
\begin{table}[ht]
\centering
  \centering
  \caption*{Table~1\hskip1em This is an example for manually numbered table, which
    would not appear in the list of tables}
  \label{tab:badtabular2}
  \begin{tabular}[c]{|c|m{0.8in}|c|c|c|c|c|}\hline
    \multicolumn{2}{|c|}{Network Topology} & \# of nodes &
    \multicolumn{3}{c|}{\# of clients} & Server \\\hline
    GT-ITM & Waxman Transit-Stub & 600 &
    \multirow{2}{2em}{2\%}&
    \multirow{2}{2em}{10\%}&
    \multirow{2}{2em}{50\%}&
    \multirow{2}{1.2in}{Max. Connectivity}\\\cline{1-3}
    \multicolumn{2}{|c|}{Inet-2.1} & 6000 & & & &\\\hline
    \multirow{2}{1in}{Xue} & Rui  & Ni &\multicolumn{4}{c|}{\multirow{2}*{\tongjithesis}}\\\cline{2-3}
    & \multicolumn{2}{c|}{ABCDEF} &\multicolumn{4}{c|}{} \\\hline
\end{tabular}
\end{table}

Roughly speaking, the simplex algorithm examines only the extreme points of the
feasible set, rather than all feasible points. At first, the simplex algorithm
selects an extreme point as the initial point. The successive extreme point is
selected so as to improve the objective function value. The procedure is
repeated until no improvement in objective function value can be made. The last
extreme point is the optimal solution.

\subsection{Nonlinear Programming}

If at least one of the functions $f(x),g_j(x),j=1,2,\cdots,p$ is nonlinear, then
SOP is called a {\em nonlinear programming}.

A large number of classical optimization methods have been developed to treat
special-structural nonlinear programming based on the mathematical theory
concerned with analyzing the structure of problems.
\begin{figure}[h]
  \centering
  \includegraphics[clip]{tongji-lib-logo.jpg}
  \caption*{Figure~1\hskip1em This is an example for manually numbered figure,
    which would not appear in the list of figures}
  \label{tab:badfigure2}
\end{figure}

Now we consider a nonlinear programming which is confronted solely with
maximizing a real-valued function with domain $\Re^n$.  Whether derivatives are
available or not, the usual strategy is first to select a point in $\Re^n$ which
is thought to be the most likely place where the maximum exists. If there is no
information available on which to base such a selection, a point is chosen at
random. From this first point an attempt is made to construct a sequence of
points, each of which yields an improved objective function value over its
predecessor. The next point to be added to the sequence is chosen by analyzing
the behavior of the function at the previous points. This construction continues
until some termination criterion is met. Methods based upon this strategy are
called {\em ascent methods}, which can be classified as {\em direct methods},
{\em gradient methods}, and {\em Hessian methods} according to the information
about the behavior of objective function $f$. Direct methods require only that
the function can be evaluated at each point. Gradient methods require the
evaluation of first derivatives of $f$. Hessian methods require the evaluation
of second derivatives. In fact, there is no superior method for all
problems. The efficiency of a method is very much dependent upon the objective
function.

\subsection{Integer Programming}

{\em Integer programming} is a special mathematical programming in which all of
the variables are assumed to be only integer values. When there are not only
integer variables but also conventional continuous variables, we call it {\em
  mixed integer programming}. If all the variables are assumed either 0 or 1,
then the problem is termed a {\em zero-one programming}. Although integer
programming can be solved by an {\em exhaustive enumeration} theoretically, it
is impractical to solve realistically sized integer programming problems. The
most successful algorithm so far found to solve integer programming is called
the {\em branch-and-bound enumeration} developed by Balas (1965) and Dakin
(1965). The other technique to integer programming is the {\em cutting plane
  method} developed by Gomory (1959).

\hfill\textit{Uncertain Programming\/}\quad(\textsl{BaoDing Liu, 2006.2})

\section*{References}
\noindent{\itshape NOTE: these references are only for demonstration, they are
  not real citations in the original text.}

\begin{enumerate}[{$[$}1{$]$}]
\item Donald E. Knuth. The \TeX book. Addison-Wesley, 1984. ISBN: 0-201-13448-9
\item Paul W. Abrahams, Karl Berry and Kathryn A. Hargreaves. \TeX\ for the
  Impatient. Addison-Wesley, 1990. ISBN: 0-201-51375-7
\item David Salomon. The advanced \TeX book.  New York : Springer, 1995. ISBN:0-387-94556-3
\end{enumerate}

\chapter{外文资料的调研阅读报告或书面翻译}
\section{单目标规划}
北冥有鱼，其名为鲲。鲲之大，不知其几千里也。化而为鸟，其名为鹏。鹏之背，不知其几
千里也。怒而飞，其翼若垂天之云。是鸟也，海运则将徙于南冥。南冥者，天池也。
\begin{equation}\tag*{(123)}
 p(y|\mathbf{x}) = \frac{p(\mathbf{x},y)}{p(\mathbf{x})}=
\frac{p(\mathbf{x}|y)p(y)}{p(\mathbf{x})}
\end{equation}

吾生也有涯，而知也无涯。以有涯随无涯，殆已！已而为知者，殆而已矣！为善无近名，为
恶无近刑，缘督以为经，可以保身，可以全生，可以养亲，可以尽年。

\subsection{线性规划}
庖丁为文惠君解牛，手之所触，肩之所倚，足之所履，膝之所倚，砉然响然，奏刀騞然，莫
不中音，合于桑林之舞，乃中经首之会。
\begin{table}[ht]
\centering
  \centering
  \caption*{表~1\hskip1em 这是手动编号但不出现在索引中的一个表格例子}
  \label{tab:badtabular3}
  \begin{tabular}[c]{|c|m{0.8in}|c|c|c|c|c|}\hline
    \multicolumn{2}{|c|}{Network Topology} & \# of nodes &
    \multicolumn{3}{c|}{\# of clients} & Server \\\hline
    GT-ITM & Waxman Transit-Stub & 600 &
    \multirow{2}{2em}{2\%}&
    \multirow{2}{2em}{10\%}&
    \multirow{2}{2em}{50\%}&
    \multirow{2}{1.2in}{Max. Connectivity}\\\cline{1-3}
    \multicolumn{2}{|c|}{Inet-2.1} & 6000 & & & &\\\hline
    \multirow{2}{1in}{Xue} & Rui  & Ni &\multicolumn{4}{c|}{\multirow{2}*{\tongjithesis}}\\\cline{2-3}
    & \multicolumn{2}{c|}{ABCDEF} &\multicolumn{4}{c|}{} \\\hline
\end{tabular}
\end{table}

文惠君曰：“嘻，善哉！技盖至此乎？”庖丁释刀对曰：“臣之所好者道也，进乎技矣。始臣之
解牛之时，所见无非全牛者；三年之后，未尝见全牛也；方今之时，臣以神遇而不以目视，
官知止而神欲行。依乎天理，批大郤，导大窾，因其固然。技经肯綮之未尝，而况大坬乎！
良庖岁更刀，割也；族庖月更刀，折也；今臣之刀十九年矣，所解数千牛矣，而刀刃若新发
于硎。彼节者有间而刀刃者无厚，以无厚入有间，恢恢乎其于游刃必有余地矣。是以十九年
而刀刃若新发于硎。虽然，每至于族，吾见其难为，怵然为戒，视为止，行为迟，动刀甚微，
謋然已解，如土委地。提刀而立，为之而四顾，为之踌躇满志，善刀而藏之。”

文惠君曰：“善哉！吾闻庖丁之言，得养生焉。”


\subsection{非线性规划}
孔子与柳下季为友，柳下季之弟名曰盗跖。盗跖从卒九千人，横行天下，侵暴诸侯。穴室枢
户，驱人牛马，取人妇女。贪得忘亲，不顾父母兄弟，不祭先祖。所过之邑，大国守城，小
国入保，万民苦之。孔子谓柳下季曰：“夫为人父者，必能诏其子；为人兄者，必能教其弟。
若父不能诏其子，兄不能教其弟，则无贵父子兄弟之亲矣。今先生，世之才士也，弟为盗
跖，为天下害，而弗能教也，丘窃为先生羞之。丘请为先生往说之。”
\begin{figure}[h]
  \centering
  \includegraphics{hello}
  \caption*{图~1\hskip1em 这是手动编号但不出现索引中的图片的例子}
  \label{tab:badfigure3}
\end{figure}

柳下季曰：“先生言为人父者必能诏其子，为人兄者必能教其弟，若子不听父之诏，弟不受
兄之教，虽今先生之辩，将奈之何哉？且跖之为人也，心如涌泉，意如飘风，强足以距敌，
辩足以饰非。顺其心则喜，逆其心则怒，易辱人以言。先生必无往。”

孔子不听，颜回为驭，子贡为右，往见盗跖。

\subsection{整数规划}
盗跖乃方休卒徒大山之阳，脍人肝而餔之。孔子下车而前，见谒者曰：“鲁人孔丘，闻将军
高义，敬再拜谒者。”谒者入通。盗跖闻之大怒，目如明星，发上指冠，曰：“此夫鲁国之
巧伪人孔丘非邪？为我告之：尔作言造语，妄称文、武，冠枝木之冠，带死牛之胁，多辞缪
说，不耕而食，不织而衣，摇唇鼓舌，擅生是非，以迷天下之主，使天下学士不反其本，妄
作孝弟，而侥幸于封侯富贵者也。子之罪大极重，疾走归！不然，我将以子肝益昼餔之膳。”


\chapter{其它附录}
其它附录的内容可以放到这里，当然如果你愿意，可以把这部分也放到独立的文件中，然后
将其\verb|\input| 到主文件中。
\end{appendix}

% 个人简历
\begin{resume}
%\resumeitem{个人简历:}
%\noindent xxxx 年 xx 月 xx 日出生于 xx 省 xx 县。\\
%\noindent xxxx 年 9 月考入 xx 大学 xx 系 xx 专业，xxxx 年 7 月本科毕业并获得 xx 学士学位。\\
%\noindent xxxx 年 9 月免试进入 xx 大学 xx 系攻读 xx 学位至今。

\resumeitem{发表论文:} % 发表的和录用的合在一起
\begin{enumerate}[{[}1{]}]
\item Improved Rule Installation for Real-Time Query Service in Software-Defined Internet of Vehicles, IEEE Transactions on Intelligent Transportation Systems, 2017. (第一作者)
\item Toward the First SDN Programming Capacity Theorem on Realizing High-Level Programs on Low-Level Datapaths, INFOCOM 2018. (共同一作)
\item Resolving the Loop in High-Level SDN Program for Multi-Table Pipeline Compilation, SmartCom 2019. (第一作者，已接收)
\item Dandelion: A Novel, High-Level Programming System for Software Defined Coalitions with Local State Sharing, MILCOM 2019. (第一作者，已接收)
%\item Yang Y, Ren T L, Zhang L T, et al. Miniature microphone with silicon-
%  based ferroelectric thin films. Integrated Ferroelectrics, 2003,
%  52:229-235. (SCI 收录, 检索号:758FZ.)
%\item 杨轶, 张宁欣, 任天令, 等. 硅基铁电微声学器件中薄膜残余应力的研究. 中国机
%  械工程, 2005, 16(14):1289-1291. (EI 收录, 检索号:0534931 2907.)
%\item 杨轶, 张宁欣, 任天令, 等. 集成铁电器件中的关键工艺研究. 仪器仪表学报,
%  2003, 24(S4):192-193. (EI 源刊.)
%\item Yang Y, Ren T L, Zhu Y P, et al. PMUTs for handwriting recognition. In
%  press. (已被 Integrated Ferroelectrics 录用. SCI 源刊.)
%\item Wu X M, Yang Y, Cai J, et al. Measurements of ferroelectric MEMS
%  microphones. Integrated Ferroelectrics, 2005, 69:417-429. (SCI 收录, 检索号
%  :896KM.)
%\item 贾泽, 杨轶, 陈兢, 等. 用于压电和电容微麦克风的体硅腐蚀相关研究. 压电与声
%  光, 2006, 28(1):117-119. (EI 收录, 检索号:06129773469.)
%\item 伍晓明, 杨轶, 张宁欣, 等. 基于MEMS技术的集成铁电硅微麦克风. 中国集成电路, 
%  2003, 53:59-61.
\end{enumerate}

%\resumeitem{研究成果:} % 有就写，没有就删除
%\begin{enumerate}[{[}1{]}]
%\item 任天令, 杨轶, 朱一平, 等. 硅基铁电微声学传感器畴极化区域控制和电极连接的
%  方法: 中国, CN1602118A. (中国专利公开号.)
%\item Ren T L, Yang Y, Zhu Y P, et al. Piezoelectric micro acoustic sensor
%  based on ferroelectric materials: USA, No.11/215, 102. (美国发明专利申请号.)
%\end{enumerate}

\end{resume}

\end{document}
